\chapter{Base modules}
\label{cap:base}

This chapter documents the modules that do not belong to a specific driver:
\code{main.rs}, \code{arch.rs}, and \code{sync.rs}. They provide the entry
point, the printing macros, I/O port access,
shutdown/reboot, and the console spinlock.

\section{\code{main.rs}: Rust entry point}

\subsection{Global attributes}

\begin{lstlisting}[caption={Attributes and modules declared in \code{main.rs}.},label={lst:mainattrs}]
#![no_std]
#![no_main]

mod arch;
mod boot;
mod gdt;
mod idt;
mod keyboard;
mod paging;
mod serial;
mod shell;
mod sync;
mod user;
mod vga;
\end{lstlisting}

\begin{itemize}
  \item \code{no\_std}: the crate does not use the standard library; there is no
        \code{std::io}, \code{String}, \code{Vec}, filesystem, or
        \emph{threads}.
  \item \code{no\_main}: the Rust runtime \code{main} function does not exist.
        The real entry point is \code{\_start}, defined in the
        assembly of \code{boot.rs}.
  \item \code{mod boot}: required for the linker to include the assembly code
        and the Multiboot2 header.
\end{itemize}

\subsection{\code{kernel\_main}}

\begin{lstlisting}[caption={Signature and main body of \code{kernel\_main}.},label={lst:kernelmain}]
#[unsafe(no_mangle)]
pub extern "C" fn kernel_main(multiboot_magic: u32, multiboot_info: u32) -> ! {
    vga::init();
    serial::init();
    paging::init(multiboot_magic, multiboot_info);
    gdt::init();
    keyboard::init();
    idt::init();
    arch::enable_interrupts();

    println!();
    println!("==============================================================");
    println!("  Rust VGA Kernel - x86_64 / modo longo");
    /* ... banner ... */

    if multiboot_magic != 0x36d7_6289 {
        println!("AVISO: assinatura multiboot2 invalida: {:#010x}",
                 multiboot_magic);
    }
    println!("Multiboot2: magic={:#010x} info={:#010x}",
             multiboot_magic, multiboot_info);
    println!();
    println!("Digite 'help' para ver os comandos.");
    println!();

    shell::run()
}
\end{lstlisting}

\begin{table}[H]
\centering
\caption{Contract of \code{kernel\_main}.}
\label{tab:kernelmain}
\begin{tabularx}{\linewidth}{@{}lX@{}}
\toprule
\textbf{Aspect} & \textbf{Description} \\
\midrule
Visibility & \code{pub extern "C"}, exported without \emph{mangling}. \\
Arguments & \code{multiboot\_magic} (must be \code{0x36D76289}) and
\code{multiboot\_info} (physical address of the Multiboot2 structure). \\
Return & \code{!}: never returns; ends in \code{shell::run()}. \\
Preconditions & CPU in long mode, paging active, valid higher-half
stack, \code{.bss} zeroed by the loader. \\
Effects & Initializes VGA, serial, frame allocator, GDT/TSS, keyboard,
IDT/PIC/PIT, and enables interrupts; prints the banner; enters the shell. \\
Called by & Assembly stub in \code{boot.rs}. \\
\bottomrule
\end{tabularx}
\end{table}

The \code{multiboot\_info} value is used by \code{paging::init} to
walk the Multiboot2 tags and free the available RAM frames. The
structure is also reserved so that it is not handed out by the allocator.

\subsection{\code{print!} and \code{println!} macros}

The macros are defined with \code{macro\_rules!} and \code{\$crate}, so that
they work in any module of the crate without imports:

\begin{lstlisting}[caption={Printing macros without allocation.},label={lst:macros}]
#[macro_export]
macro_rules! print {
    ($($arg:tt)*) => {
        $crate::vga::_print(format_args!($($arg)*))
    };
}

#[macro_export]
macro_rules! println {
    () => { $crate::print!("\n") };
    ($($arg:tt)*) => {
        $crate::print!("{}\n", format_args!($($arg)*))
    };
}
\end{lstlisting}

\code{format\_args!} creates a formatting description without allocating;
the \code{fmt::Write} implementation in \code{vga.rs} consumes the
fragments. There is no \code{format!} (which would require \code{alloc}).

\subsection{Panic handler}

\begin{lstlisting}[caption={\emph{Panic handler} of the kernel.},label={lst:panic}]
#[panic_handler]
fn panic(info: &PanicInfo) -> ! {
    crate::println!();
    crate::println!("[PANIC] {}", info);
    crate::arch::halt_loop()
}
\end{lstlisting}

Because the profile uses \code{panic = "abort"}, the handler does not need
\emph{unwinding}. It prints the message to the console and halts the CPU. If the
panic occurs before \code{vga::init()}, the output may not be
visible; if it occurs inside the console with the spinlock held, a
\emph{deadlock} may occur (the handler would try to print and would wait for the lock).

\section{\code{arch.rs}: architecture and I/O ports}

\subsection{I/O ports}

The \code{arch::port} module encapsulates the x86 \code{in} and
\code{out} instructions:

\begin{lstlisting}[caption={Reading and writing 8-bit ports.},label={lst:portio}]
pub unsafe fn inb(port: u16) -> u8 {
    let value: u8;
    asm!("in al, dx",
         in("dx") port,
         out("al") value,
         options(nomem, nostack, preserves_flags));
    value
}

pub unsafe fn outb(port: u16, value: u8) {
    asm!("out dx, al",
         in("dx") port,
         in("al") value,
         options(nomem, nostack, preserves_flags));
}
\end{lstlisting}

The functions are \code{unsafe} because any port may not exist
or may have side effects on the hardware. The \code{nomem},
\code{nostack}, and \code{preserves\_flags} options tell the compiler that the
instruction does not change memory, stack, or flags, allowing optimizations
and preventing the block from being removed.

\begin{table}[H]
\centering
\caption{API of \code{arch}.}
\label{tab:archapi}
\begin{tabularx}{\linewidth}{@{}lX@{}}
\toprule
\textbf{Function} & \textbf{Description} \\
\midrule
\code{inb(port)} & Reads a byte from an I/O port. \code{unsafe}. \\
\code{outb(port, value)} & Writes a byte to a port. \code{unsafe}. \\
\code{halt\_loop()} & Executes \code{cli; hlt} in an infinite loop; does not
return. Used in the panic handler and in the \code{halt} command. \\
\code{reboot()} & Waits for the 8042 input buffer to empty and sends the
\code{0xFE} command (reset). If it does not work, it falls back to \code{halt\_loop}. \\
\bottomrule
\end{tabularx}
\end{table}

\subsection{Reboot}

The reset uses the 8042 controller, which historically controls the CPU reset
line. The code waits up to 100,000 iterations for the
\code{INPUT\_FULL} bit (\code{0x02}) of the status register to clear and writes \code{0xFE} to the
command port (\code{0x64}). On virtualized machines or machines without an 8042, the
reset may not occur; in that case the CPU is halted.

\section{\code{sync.rs}: spinlock}

\subsection{Implementation}

\begin{lstlisting}[caption={Spinlock with \code{AtomicBool}.},label={lst:spinlock}]
pub struct SpinLock<T> {
    locked: AtomicBool,
    data: UnsafeCell<T>,
}

unsafe impl<T: Send> Sync for SpinLock<T> {}
unsafe impl<T: Send> Send for SpinLock<T> {}

impl<T> SpinLock<T> {
    pub const fn new(data: T) -> Self { /* ... */ }

    pub fn lock(&self) -> SpinLockGuard<'_, T> {
        while self.locked.swap(true, Ordering::Acquire) {
            core::hint::spin_loop();
        }
        SpinLockGuard { lock: self }
    }
}
\end{lstlisting}

The \emph{guard} implements \code{Deref}, \code{DerefMut}, and \code{Drop}. When it is
destroyed, it stores \code{false} with \code{Release} ordering. The
spinlock is used only for the console:

\begin{lstlisting}[caption={Global console protected by a spinlock.},label={lst:consoleglobal}]
static CONSOLE: SpinLock<Console> = SpinLock::new(Console::new());
\end{lstlisting}

Because \code{Console::new()} is \code{const}, the lock can be initialized
statically without runtime initialization code.

\subsection{Properties and limitations}

\begin{itemize}
  \item \textbf{Spinning lock}: if another CPU or an interrupt
        holds the lock, the current CPU spins until it is released. There is no
        \emph{sleep} or timeout.
  \item \textbf{No IRQ save/restore}: disabling interrupts when
        acquiring the lock is not done. With IRQs disabled at boot, this is
        not a problem today; once IRQs are enabled, a handler that uses the
        console may cause a \emph{deadlock}.
  \item \textbf{No recursion detection}: acquiring the same lock twice
        on the same CPU hangs.
  \item \textbf{Memory ordering}: \code{Acquire}/\code{Release} guarantees
        that writes to the console are not reordered out of the
        critical section on architectures with weak memory; on x86,
        the model is stronger, but the annotation keeps the code
        correct.
\end{itemize}

\begin{nota}[Current scope]
With one CPU, interrupts disabled, and only the shell calling
\code{print!}, the spinlock is never contended. It exists as preparation
for extensions and to document the intent of mutual exclusion.
\end{nota}
